\section{Economic Realization}

We finally examine how residual signals are realized in future returns. The goal is not to prove a single causal mechanism, but to characterize horizon and state dependence. The first set of tables concerns the three-model CNN--MLP-4--HistGB decomposition; the final subsection separately reports pairwise temporal CNN--1D-CNN residuals.

\begin{table}[htbp]
\centering
\caption{Economic realization of model-specific signals}
\label{tab:mechanism}
\small
\begin{tabular}{p{0.23\linewidth}p{0.32\linewidth}p{0.32\linewidth}}
\toprule
Dimension & CNN-specific & MLP-specific \\
\midrule
Horizon & Short-lived; strongest at $t+1$ & Persistent; strong through days $2$--$5$ \\
Session & Intraday-leaning, weak evidence & No clear ON/ID split \\
State dependence & Turnover / pressure amplified & Turnover, momentum, volatility amplified \\
Pressure interaction & Positive FM interaction & Approximately zero interaction \\
Direction & Stronger after negative recent returns & Stronger after positive recent returns over $R_{1:5}$ \\
Institutional states & Suggestive, not causal & Not T+1-pressure centered \\
\bottomrule
\end{tabular}
\begin{minipage}{0.92\linewidth}
\footnotesize
\textit{Notes:} This table summarizes evidence from horizon decomposition, overnight/intraday tests, pressure interactions, continuation-versus-reversal sorts, and T+1 / price-limit screening tests.
\end{minipage}
\end{table}

\begin{table}[htbp]
\centering
\caption{Pressure-state Rank IC}
\label{tab:mechanism_pressure}
\begin{tabular}{llrrrrr}
\toprule
Return & Score & Low & Mid & High & High--Low & NW $t$ \\
\midrule
$R_{1:5}$ & CNN-specific & 0.0133 & 0.0141 & 0.0172 & 0.0040 & 1.38 \\
$R_{1:5}$ & MLP-4-specific & 0.0253 & 0.0315 & 0.0458 & 0.0205 & 7.36 \\
$r_1$ & CNN-specific & 0.0161 & 0.0147 & 0.0216 & 0.0055 & 1.97 \\
$r_1$ & MLP-4-specific & 0.0100 & 0.0169 & 0.0202 & 0.0103 & 3.68 \\
\bottomrule
\end{tabular}
\begin{minipage}{0.92\linewidth}
\footnotesize
\textit{Notes:} Pressure is a pre-specified composite of turnover proxy, drawdown pressure, and terminal volume intensity. Stocks are sorted into daily pressure terciles.
\end{minipage}
\end{table}

\begin{table}[htbp]
\centering
\caption{Fama--MacBeth pressure interactions}
\label{tab:mechanism_interaction}
\begin{tabular}{llrr}
\toprule
Return & Score & Main score $\bar{\beta}_1$ ($t$) & Interaction $\bar{\beta}_3$ ($t$) \\
\midrule
$R_{1:5}$ & CNN-specific & 0.00043 (4.22) & 0.00035 (4.26) \\
$R_{1:5}$ & MLP-4-specific & 0.00124 (10.41) & 0.00004 (0.54) \\
$r_1$ & CNN-specific & 0.00034 (7.28) & 0.00024 (7.53) \\
$r_1$ & MLP-4-specific & 0.00031 (6.72) & 0.00000 (0.01) \\
\bottomrule
\end{tabular}
\begin{minipage}{0.92\linewidth}
\footnotesize
\textit{Notes:} Regressions include score, pressure, score-by-pressure interaction, size, and recent return controls. Newey--West lag is 4.
\end{minipage}
\end{table}

\begin{table}[htbp]
\centering
\caption{Continuation versus rebound: high-minus-low specific-score spreads}
\label{tab:mechanism_direction}
\begin{tabular}{llrrrr}
\toprule
Return & Score & Negative recent & $t$ & Positive recent & $t$ \\
\midrule
$R_{1:5}$ & CNN-specific & 0.00176 & 5.59 & 0.00076 & 1.98 \\
$R_{1:5}$ & MLP-4-specific & 0.00201 & 6.42 & 0.00379 & 9.66 \\
$r_1$ & CNN-specific & 0.00087 & 6.60 & 0.00072 & 4.15 \\
$r_1$ & MLP-4-specific & 0.00074 & 5.50 & 0.00014 & 0.78 \\
\bottomrule
\end{tabular}
\begin{minipage}{0.92\linewidth}
\footnotesize
\textit{Notes:} Stocks are first grouped by recent cumulative return terciles. Within negative and positive recent-return groups, the table reports future-return spreads between high and low specific-score portfolios.
\end{minipage}
\end{table}


In the three-model decomposition, CNN-specific predictability is short-lived. It is strongest at the next-day horizon and decays quickly. It is also amplified in pressure-like states, especially when turnover and recent-path stress are high. In Fama--MacBeth specifications, interactions between CNN-specific scores and pressure states are positive and significant. The signal is more rebound-leaning or mixed than pure continuation.

In the three-model decomposition, MLP-4-specific predictability is more persistent. It remains strong through days two to five and is amplified by turnover, recent momentum, and volatility. Its five-day spread is larger following positive recent returns, consistent with continuation. T+1-adjacent and price-limit state tests are suggestive rather than causal: they are useful for screening whether institutional frictions are state-consistent with the signal, but they do not prove that such frictions generate the alpha.

Within the three-model decomposition, the economic interpretation is therefore asymmetric. The CNN-specific signal resembles a short-window path-pressure signal with rebound-like elements, while the MLP-4-specific signal resembles a more persistent continuation signal. This distinction supports the broader claim that different nonlinear model classes extract different economically relevant structures from the same price history.

\subsection{CNN versus Temporal CNN}

Table~\ref{tab:temporal_mechanism} provides the parallel E5T test for \emph{pairwise temporal} CNN--1D-CNN residuals. The CNN-specific component conditional on 1D-CNN is concentrated at the next-day horizon and decays over longer windows. In contrast, the 1D-CNN-specific component conditional on CNN increases from an $r_1$ Rank IC of 0.0137 to 0.0407 over $R_{1:5}$ and 0.0428 over $R_{1:10}$. The latter residual is substantially stronger in high-pressure states and remains positive in both negative and positive recent-return groups. The pairwise CNN-specific residual has much smaller medium-horizon spreads and is relatively stronger after negative recent returns. Price-limit and T+1-state results remain exploratory because such states are sparse.

\begin{table}[htbp]
\centering
\caption{Economic realization: pairwise CNN--1D-CNN specific signals}
\label{tab:temporal_mechanism}
\resizebox{\textwidth}{!}{%
\begin{tabular}{lrrrrrr}
\toprule
Component & $r_1$ IC & $R_{1:5}$ IC & $R_{1:10}$ IC & Pressure $\Delta$IC & Neg. recent spread & Pos. recent spread \\
\midrule
CNN-specific conditional on 1D-CNN & 0.0184 & 0.0140 & 0.0089 & 0.0007 (0.25) & 0.00125 & 0.00012 \\
1D-CNN-specific conditional on CNN & 0.0137 & 0.0407 & 0.0428 & 0.0250 (7.39) & 0.00429 & 0.00447 \\
\bottomrule
\end{tabular}%
}
\vspace{0.25em}
\footnotesize\textit{Notes:} IC is mean daily Rank IC. Pressure $\Delta$IC is high-minus-low pressure-state $R_{1:5}$ IC, with Newey--West $t$ in parentheses. Spreads are $R_{1:5}$ high-minus-low specific-score spreads within CumRet5 terciles.
\end{table}

